%%% Формат и шрифты
\documentclass[a4paper]{article}
\usepackage{setspace} 
\onehalfspacing					% двойной интервал 
\usepackage[14pt]{extsizes}		% Увеличенный 14 шрифт
\usepackage{cmap}				% Улучшенный поиск русских слов в полученном pdf-файле
\usepackage[T2A]{fontenc}		% Поддержка русских букв
\usepackage[utf8]{inputenc}		% Кодировка utf8
\usepackage[english, russian]{babel}
\renewcommand{\rmdefault}{cmr}
% Шрифт -- Computer Modern Roman
% Языки: русский, английский
\usepackage{color}				%Для работы с цветами
\usepackage{hyperref}
\hypersetup{pdfborder={0 0 0}, colorlinks = true, citecolor = blue, linkcolor = black}
% Гиперссылки в документе кликаются, разделы в оглавлении тоже
% Для красоты ссылки на литературу подкрашиваются синим. Если не хотите, можно это и убрать, заменив blue на black
\usepackage{comment}			%Для комментирования больших участков кода
\usepackage{lipsum} 			%Если надо заполнить страницу шаблонным текстом

%%% Выравнивание и переносы %%%
\sloppy					% Избавляемся от переполнений
\clubpenalty=10000		% Запрещаем разрыв страницы после первой строки абзаца
\widowpenalty=10000		% Запрещаем разрыв страницы после последней строки абзаца
\parindent=10mm
\usepackage{indentfirst} % отделять первую строку раздела абзацным отступом тоже
\usepackage[left=30mm, top=15mm, right=15mm, bottom=20mm, nohead, footskip=8mm]{geometry} % настройки полей документа


%%% Математические пакеты %%%
\usepackage{amsthm,amsfonts,amsmath,amssymb,amscd,mathtools}
\usepackage{algorithm}
\usepackage{algpseudocode}

\usepackage[most]{tcolorbox}

%%% Разделы %%%
% Делаем названия большими жирными буквами
\usepackage{titlesec}
\titleformat{\section}
{\centering\large\bfseries}
{\thesection. }
{0em}{}
\titlespacing*{\section}{\parindent}{*2}{*2}% Настройка вертикальных и горизонтальных отступов
\titleformat{\subsection}
{\centering\large\bfseries}
{\thesubsection. }
{0em}{}
\titlespacing*{\subsection}{\parindent}{*2}{*2}% Настройка вертикальных и горизонтальных отступов
\titleformat{\subsubsection}
{\centering\large\bfseries}
{\thesubsubsection. }
{0em}{}
\titlespacing*{\subsubsection}{\parindent}{*2}{*2}% Настройка вертикальных и горизонтальных отступов

%%% БИБЛИОГРАФИЯ %%% 
%для сокращения расстояния между элементами библиографии используем 
\let\oldthebibliography=\thebibliography 
\let\endoldthebibliography=\endthebibliography 
\renewenvironment{thebibliography}[1] 
{ \begin{oldthebibliography}{#1} \setlength{\parskip}{0.5ex} 
		\setlength{\itemsep}{-0.5ex} } 
	{\end{oldthebibliography}} 
\bibliographystyle{utf8gost705u} % стилевой файл для оформления по ГОСТу 
\makeatletter 
\renewcommand{\@biblabel}[1]{#1.} % Заменяем библиографию с квадратных скобок на точку 
\makeatother 

%%% Нумерация страниц %%%
\usepackage{fancyhdr}
\pagestyle{fancy}
\fancyhf{}
\fancyfoot[C]{\normalsize\thepage}				%Подпись внизу страницы по центру, номер 14-м шрифтом
\fancyheadoffset{0mm}
\fancyfootoffset{0mm}
\setlength{\headheight}{17pt}
\renewcommand{\headrulewidth}{0pt}
\renewcommand{\footrulewidth}{0pt}
\fancypagestyle{plain}{ 
	\fancyhf{}
	\rhead{\thepage}}
\setcounter{page}{1} % начать нумерацию страниц с №1

%%%Таблицы, картинки и листинги кода%%%
\usepackage{array}					%Пакет для написания матриц и массивов
%\usepackage{tabu}					%Продвинутый пакет для таблиц
\usepackage{multirow}				%Для объединения ячеек в строке таблицы
\usepackage{multicol}				%Для объединения ячеек в столбце таблицы
\usepackage{float}					%Для работы с плавающими объектами: картинками, таблицами, листингами и т.д.
\usepackage{graphicx}				%Для изображений
\usepackage{subfigure}				%Для нумерации отдельных картинок внутри составной
\usepackage{wrapfig}				%Для картинок внутри текста
\usepackage[nooneline]{caption}
\captionsetup[table]{
	skip=0pt,
	singlelinecheck=off} 			%Подпись к таблице сверху слева
\captionsetup[figure]{
	justification=centering,
	labelsep=endash} 				%Подпись к картинке по центру внизу
\usepackage{listings}				%Для написания листингов кода. Цвета подобраны под язык С++
\definecolor{mygray}{rgb}{0.4,0.4,0.4}
\definecolor{mygreen}{rgb}{0,0.8,0.6}
\definecolor{myorange}{rgb}{1.0,0.4,0}
\lstset{
	language=C++, 
	commentstyle=\color{mygray},
	numberstyle=\tiny\color{mygray},
	keywordstyle=\color{cyan},
	stringstyle=\color{myorange},
	basicstyle=\footnotesize\ttfamily\color{black},
	numbers=left,
	%frame=single					%Если нужна рамка, можно раскомментировать
}

%%% Определение теорем, лемм, утверждений и т.д. %%%
\newtheorem{theorem}{Теорема}
\newtheorem{definition}{Определение}
\newtheorem{lemma}{Лемма}
\newtheorem{proposition}{Утверждение}
\newtheorem{corollary}{Следствие}
%\renewcommand{\qedsymbol}{$\blacksquare$}%Если хочется черный квадрат вместо пустого в конце доказательства.
\addto\captionsrussian{\def\refname{СПИСОК ЛИТЕРАТУРЫ}}


\begin{document}
	\begin{singlespacing} % На титульнике одинарный интервал
		\thispagestyle{empty} % Не нумеруем титульный лист
		\begin{center}
			\small{МИНИСТЕРСТВО НАУКИ И ВЫСШЕГО ОБРАЗОВАНИЯ РОССИЙСКОЙ ФЕДЕРАЦИИ}\\
			\bigskip
			\small{ФЕДЕРАЛЬНОЕ ГОСУДАРСТВЕННОЕ АВТОНОМНОЕ ОБРАЗОВАТЕЛЬНОЕ УЧРЕЖДЕНИЕ ВЫСШЕГО ОБРАЗОВАНИЯ}\\
			\medskip
			\small{«НОВОСИБИРСКИЙ НАЦИОНАЛЬНЫЙ ИССЛЕДОВАТЕЛЬСКИЙ ГОСУДАРСТВЕННЫЙ УНИВЕРСИТЕТ» (НОВОСИБИРСКИЙ ГОСУДАРСТВЕННЫЙ УНИВЕРСИТЕТ, НГУ)}\\ 
			\hfill \break
			\small{Механико-математический факультет}\\
			\medskip
			\small{Кафедра теоретической кибернетики}\\
			\smallskip
			\bigskip
			\small{Направление подготовки «Прикладная математика и информатика», бакалавриат}\\
			\vspace{2.5cm}
			\small{\textbf{ВЫПУСКНАЯ КВАЛИФИКАЦИОННАЯ РАБОТА БАКАЛАВРА}}\\
			\medskip
			\hfill \break
			\small{Белькова Татьяна Дмитриевна}\\
			\hfill \break
			\normalsize{\textbf{Матэвристика для задачи упаковки нерегулярных многоугольников в рюкзак}}\\
			%\hfill \break
		\end{center}
	\end{singlespacing}
	\vspace{2.5cm}
	\noindent
	\begin{minipage}{70mm}\parindent=0em
	\begin{flushleft}
			\small{\textbf{«К защите допущена»}\par
			\medskip
			Заведующий кафедрой\par
			\medskip
			*ученая степень*, *ученое звание*,\par
			\medskip		
			\small{\makebox[3.5cm]{Тахонов И.И.}\;\;}% Для выравнивания придется поиграться чиселками
			/ \makebox[2cm]{\dotfill}\par}
		\scriptsize{~(Фамилия, И.О.)~/~(подпись, МП)}\par
		\smallskip
		\small{«\makebox[1cm]{\dotfill}»\makebox[3cm]{\dotfill} 2026г.}
	\end{flushleft}
	\end{minipage}
	\hspace{3cm}
	\begin{minipage}{80mm}\parindent=0em
		\begin{flushleft}
			\small{\textbf{Научный руководитель}\par
				\medskip
				*Должность* *Аббревиатура кафедры*\par
				\medskip
				*ученая степень*, *ученое звание*\par
				\medskip		
				\small{\makebox[3.5cm]{Шперлинг С.М.}\;}% Для выравнивания придется поиграться чиселками
				/ \makebox[2cm]{\dotfill}\par}
			\scriptsize{~(Фамилия,И.О.)~/~(подпись, МП)}\par
			\smallskip
			\small{«\makebox[1cm]{\dotfill}»\makebox[3cm]{\dotfill} 2026г.}
		\end{flushleft}
	\end{minipage}
	\vspace{1cm}
	\begin{flushright}
		\small Дата защиты: \small{«\makebox[1cm]{\dotfill}»\makebox[3cm]{\dotfill} 2026г.}
	\end{flushright}
    \vspace{1cm}
		\begin{center}
			\small{Новосибирск, 2026}
		\end{center} 

	\newpage
	\thispagestyle{empty} % Не нумеруем реферат
	\begin{center}
		\large{\textbf {Реферат}}
	\end{center} 	

    Название работы: Матэвристика для задачи упаковки нерегулярных многоугольников в рюкзак.
    
    Здесь идет текст реферата.  
    
    \textit{Количество страниц:} 12
    
    \textit{Количество использованных источников:} 4
    
    \textit{Количество иллюстраций:} 3
    
    \textit{Количество таблиц:} 1
    
    \textit{Ключевые слова:} ключевое слово 1, ключевое слово 2.
	
	%CОДЕРЖАНИЕ
	\newpage
	
	\tableofcontents
	\normalsize
	
	%ВВЕДЕНИЕ	
	\newpage
	\section*{ВВЕДЕНИЕ}
	\addcontentsline{toc}{section}{ВВЕДЕНИЕ} %Не нумеруем введение, но добавляем в оглавление

    Задача о рюкзаке это классическая задача комбинаторной оптимизации, различные постановки которой применяются в логистике, производственном планировании и задачах раскроя материалов. В представленной выпускной квалификационной работе бакалавра рассматривается задача упаковки нерегулярных многоугольников в прямоугольный рюкзак фиксированного размера.
    
    Пусть задан набор из $N$ нерегулярных многоугольников. Требуется выбрать такой поднабор объектов и определить их положение в рюкзаке, чтобы выполнялись следующие условия:

    \vspace{-0.28cm}
    \begin{itemize}
        \setlength{\itemsep}{0.15em}
        \setlength{\parskip}{0pt}
        \setlength{\parsep}{0pt}

        \item все выбранные многоугольники размещаются внутри рюкзака;
        \item выбранные многоугольники не пересекаются между собой;
        \item для каждого многоугольника допускается выбор одного из допустимых поворотов;
        \item суммарная площадь выбранных многоугольников максимальна среди всех допустимых размещений.
    \end{itemize}
    
    Задачи упаковки нерегулярных фигур актуальны на производствах, связанных с раскроем материала. Практическая значимость заключается в необходимости минимизации расходов и максимизации использованного места на полотне. Часто в рамках таких задач приходится находить равновесие между вычислительной сложностью и качеством найденного решения алгоритма. 

    Существенное влияние на вычислительную сложность алгоритма оказывают инструменты, используемые для представления расположения объектов и определения пересечения многоугольников.

    Первым из таких инструментов является способ представления расположения фигур в области упаковки. Их можно разделить на два типа: дискретные и непрерывные методы. Оба метода имеют свои недостатки: при применении дискретного метода имеются неточности в геометрии фигуры; при использовании непрерывного метода геометрия прямоугольника представлена точно, однако метод сложен вычислительно. Основная проблема использования? дискретного метода заключается в неточности представления геометрии фигуры,. Ттогда как непрерывный метод точно представляет геометрию многоугольника, но вычислительно очень сложен. В рамках данной работы выбран полу-дискретный метод представления объектов, как компромисс между вычислительной сложностью и точностью представления геометрии многоугольников. Данный метод представлен и использован как инструмент в жадном алгоритме bottom-left fill Chehrazad S., Roose D. и Wauters T \cite{BLF semi}. 

    Следующий инструмент это способ определения пересечения многоугольников при упаковке. В обзорной статье \cite{tutorial} рассмотрены основные подходы к определению пересечения нерегулярных фигур. Bennell J. A. и Oliveira J. F. выделяют raster-подходы, методы прямой тригонометрии, методы на основе no-fit polygon (no-fit многоугольник) и методы на основе phi-function. 
    
    Umetani S. и Murakami S \cite{discr strip} рассмотрели задачу упаковки в полосу с использованием raster-подхода. Авторами отмечено, что главный недостаток raster-подхода это сильная зависимость вычислительных затрат от шага сетки. 

    В работе \cite{NFP} был рассмотрен no-fit polygon как способ определения пересечений многоугольников. При этом основной практической проблемой авторы называют не саму проверку пересечения, а построение полного no-fit polygon для многоугольников.

    Сами алгоритмы решения данной задачи можно разделить на три больших группы: конструктивные эвристики, метаэвристики и математические модели. Конструктивные и метаэвристические методы позволяют быстро получать допустимые решения и улучшать их за счет перебора различных порядков размещения и локальных перестроений, однако обычно не дают гарантии близости к оптимуму. Например, в статье \cite{BLF} представлен bottom-left fill алгоритм, использующий no-fit многоугольник для определения пересечения многоугольников. Математические модели позволяют явно учитывать ограничения непересечения и размещения, но при увеличении числа объектов и росте точности дискретизации становятся вычислительно трудными. Leao A. A. S., Toledo F. M. B., Oliveira J. F., Carravilla M. A. в статье \cite{MIP strip} привели математическую модель, описывающую сходную задачу упаковки в полосу.

    Целью данной работы является разработка матэвристического подхода к задаче упаковки нерегулярных многоугольников в прямоугольный рюкзак и оценка качества полученного алгоритма. Идея алгоритма заключается в получении первичного решения с помощью bottom-left fill алгоритма \cite{BLF} и его последующем улучшении на основе нахождения доступного решения математической модели \cite{MIP strip}.

    Структура работы. to be continued...
    
    \newpage
	
	\section{Математическая постановка задачи}
	\subsection{Постановка задачи}
    Имеется прямоугольный рюкзак с известными размерами - высотой $H$ и шириной $W$, а также конечный набор из $N$ объектов. Каждый объект $i$ представлен в виде многоугольника. Задача заключается в максимизации суммарной цены многоугольников, упакованных в рюкзак. В качестве цены для каждого многоугольника $i$ возьмем его площадь $S_i$. Допускаются повороты объектов на $90^\circ$, $180^\circ$ и $270^\circ$ градусов.

    \subsection{Геометрические параметры объектов}

    Каждый объект $i, i = 1,...,N$ имеет свою точку отсчета. Точка может быть выбрана любым удобным образом. В данной работе в качестве точки отчета выбирается первая точка при объявлении многоугольника. Обозначим за $x_i^{min}$ и $x_i^{max}$ минимальное и максимальное значение координаты $x$, а $y_i^{min}$ - минимальная $y$ строка и $y_i^{max}$ - максимальная $y$ строка объекта относительно выбранной точки отсчета. Начало координат для системы координат каждого объекта располагается в его точке отсчета. 

    \begin{figure}[h!]
        \centering
        {\normalsize % Увеличение текста для всего рисунка
        \scalebox{0.7}{\input{pieceparam.pdf_tex}} }
        \caption{Пример параметров многоугольника}
		\label{Figure1}
	\end{figure}

    Для определения непересечения фигур в рамках данной работы будет применен метод no-fit polygon (no-fit многоугольник, NFP). Данный многоугольник показывает то, как мы можем расположить объект $i$ относительно объекта $j$ так, чтобы они не пересекались, но касались друг друга. Дальше в работе no-fit polygon для пары $(i, j)$ будет обозначаться как $NFP_{ij}$. Важно отметить, что $NFP_{ij}$ отличается от $NFP_{ji}$.

    \begin{figure}[h!]
    \centering
        {\huge % Увеличение текста для всего рисунка
        \scalebox{0.4}{\input{nfp_build.pdf_tex}} }
        \caption{Пример построения NFP}
		\label{Figure1}

    \end{figure}

    Представим $NFP_{ij}$ в полу-дискретной кодировке. Для каждого $NFP_{ij}$ определим $n_{ij}^{min}$ и $n_{ij}^{max}$ - минимальную и максимальную строки, на которых находится отрезки $NFP_{ij}$ в полу-дискретном представлении. В этом случае, точкой отсчета no-fit polygon является точка отчета предмета $i$. Также для каждого отрезка в полу-дискретной кодировке определим непрерывные координаты левого и правого краев отрезка $c$ на строке $k$. Обозначим их как $a_{ij}^{kc}$ и $b_{ij}^{kc}$ соответственно.

    \vspace{1cm} % 1cm пустого места
    тут будет рисуночек( 
    \vspace{1cm} % 1cm пустого места

    Так как мы рассматриваем задачу с поворотами, то нам необходимо определить дополнительные переменные для каждого объекта. Будем считать, что для каждого многоугольника $i \in \{ 1, \ldots, N\}$ определен параметр $r = 0, \ldots, R$. В общем случае угол поворота оригинального многоугольника определяется как $r \cdot \frac{360^\circ}{R}$. В данной работе рассматривается $R = 4$, то есть углы поворота определяются как $r \cdot 90^\circ$. 
    
    % \vspace{1cm} % 1cm пустого места
    % тут будет нормальное описание постановки % и что я имела тут ввиду......
    % \vspace{1cm} % 1cm пустого места

    \subsection{Переменные модели}
    
    Для определения математической модели введем следующие переменные: 
    $x_{ir}$ показывает $x$-координату (по непрерывной оси) расположения точки отсчета объекта $i$ при повороте $r$. 
    
    \[
    \delta_{ir}^{s_i} = 
     \begin{cases}
    1, & \text{если точка отсчета объекта $i$ при повороте $r$}\\ & \text{расположена на строке $s_i$},\\
    0, & \text{иначе.}
    \end{cases}
    \]

    Для $1 \leq i \leq N, 0 \leq r \leq R - 1, 0 \leq s_i \leq S - 1$

    
    Данная переменная также говорит был ли объект $i$ при повороте $r$ упакован, так как если $\sum_{s_i=0}^{S - 1} \delta_{ir}^{s_i} = 0$, то он не упакован.
    
    
    \[
    \gamma_{ir,jr'}^{c} = 
     \begin{cases}
    1, & \text{если точка отсчета объекта $i$ при повороте $r$}\\ & \text{расположена справа от объекта $j$ при повороте $r'$ на отрезке $c$},\\
    0, & \text{иначе.}
    \end{cases}
    \]
    
    Для $1 \leq i \leq j \leq N, 0 \leq r,r' \leq R - 1, c = 1,\ldots,C_{ir,jr'}^{s_j-s_i}$
    
    $s_j-s_i$ является расстоянием между объектами $i$ и $j$ при поворотах $r$ и $r'$, соответственно, по координате $y$, когда объекты расположены на строках $s_i$ и $s_j$, соответственно. $C_{ir,jr'}^{s_j-s_i}$ - количество отрезков в многоугольнике $NFP_{ir}^{jr'}$ на строке $s_j-s_i$.

	\subsection{Ограничения модели}
    
    \begin{equation}
    \sum_{s_i=0}^{S - 1 } \sum_{r=0}^{R - 1} \delta_{ir}^{s_i} \leq 1, \quad i = 1,\ldots,N
    \end{equation}
    
    \begin{equation}
    (1 - \sum_{s_i= 0 }^{S - 1} \sum_{r=0}^{R - 1} \delta_{ir}^{s_i}) \cdot M \leq s_i, \quad i = 1,\ldots,N
    \end{equation}
    
    \begin{equation}
    x_{ir} + x_{ir}^{max} \leq W + (1 - \sum_{s_i=0}^{S - 1} \sum_{r= 0}^{R - 1} \delta_{ir}^{s_i}) \cdot M, \quad i = 1,\ldots,N, \quad r = 0,\ldots,R-1
    \end{equation}
    
    \begin{equation}
    x_{ir} \geq -x_{ir}^{min}, \quad i = 1,\ldots,N, \quad r = 0,\ldots,R-1
    \end{equation}
        
    \begin{equation}
    s_i + y_i^{max} \leq H + (1 - \sum_{s_i=0}^{S - 1} \sum_{r=0}^{R - 1} \delta_{ir}^{s_i}) \cdot M, \quad i = 1,\ldots,N
    \end{equation}

    \begin{equation}
    s_i + y_i^{min} \geq 0, \quad i = 1,\ldots,N
    \end{equation}
    % в модели - s_i - координаты, вроде бы... на основе текста в строке 255 
        
    \begin{align}
    x_{ir} \leq\;& x_{jr'} - b_{ir,jr'}^{(s_j - s_i)c} + \gamma_{ir,jr'}^c \cdot M 
    + (1 - \delta_{ir}^{s_i}) \cdot M + (1 - \delta_{jr'}^{s_j}) \cdot M + \notag \\
    &+ \left(1 - \sum_{s_i=0}^{S-1} \sum_{r=0}^{R-1} \delta_{ir}^{s_i}\right) \cdot M 
    + \left(1 - \sum_{s_j=0}^{S-1} \sum_{r=0}^{R-1} \delta_{jr}^{s_j}\right) \cdot M,
    \label{eq:nonoverlap1}
    \\[0.3cm]
    x_{ir} \geq\;& x_{jr'} - a_{ir,jr'}^{(s_j - s_i)c} - (1 - \gamma_{ir,jr'}^c) \cdot M 
    - (1 - \delta_{ir}^{s_i}) \cdot M - (1 - \delta_{jr'}^{s_j}) \cdot M \notag - \\
    &- \left(1 - \sum_{s_i=0}^{S-1} \sum_{r=0}^{R-1} \delta_{ir}^{s_i}\right) \cdot M 
    - \left(1 - \sum_{s_j=0}^{S - 1} \sum_{r=0}^{R - 1} \delta_{jr}^{s_j}\right) \cdot M,
    \label{eq:nonoverlap2}
    \end{align}

    \[
    \begin{gathered}
    i,j = 1,\ldots,N, \quad i + 1 \leq j \leq N, \quad
    c = 1, \ldots,C_{ir,jr'}^{\,s_j-s_i}, \quad r,r' = 0,\ldots,R-1, \\
    n_{ir,jr'}^{\min} \leq s_j-s_i \leq n_{ir,jr'}^{\max}, \quad
    -y_i^{\min} \leq s_i \leq S - y_i^{\max}, \quad
    -y_j^{\min} \leq s_j \leq S - y_j^{\max}
    \end{gathered}
    \label{eq:nonoverlap_index}
    \]
    
    \begin{equation}
    \gamma_{ir,jr'}^c \in \{0,1\}, \hspace{0.2cm} 1 \leq i \leq j \leq N, \hspace{0.2cm} 0 \leq r, r' \leq R - 1 
    \end{equation}

    Ограничение (1) означает, что если объект упакован, то он может быть размещён только один раз и только в одном из допустимых поворотов. 
    
    Ограничение (2) задаёт условие, при котором для неупакованного предмета его полоса располагается вне рюкзака. Тем самым обеспечивается, что ограничения на непересечение учитываются только для упакованных многоугольников.
        
    Ограничение (3) гарантирует, что упакованный объект не выходит за правую границу рюкзака.
    
    Ограничение (4) устанавливает, что все многоугольники располагаются не левее левой границы рюкзака.
    
    Ограничение (5) обеспечивает выполнение условия, при котором упакованный объект не выходит за верхнюю границу рюкзака.
    
    Ограничение (6) задаёт нижнюю границу размещения и тем самым не допускает выхода многоугольников за пределы рюкзака снизу.
    
    Ограничения (7) и (8) описывают условия непересечения: для каждой пары упакованных объектов на каждой строке должно выполняться такое взаимное расположение, при котором пересечение невозможно.
    
    Ограничение (9) определяет допустимые значения переменной $\gamma_{ir,jr'}^c$.

	\subsection{Целевая функция}
    
    D качестве целевой функции выбрана максимизация суммарной площади упакованных объектов. Такой выбор соответствует практической постановке задачи, поскольку позволяет максимизировать долю использованного пространства рюкзака. В качестве ``цены'' каждого объекта используется величина $S_i$, равная площади соответствующего многоугольника.
    
    \begin{equation}
    \max \sum_{i=1}^{N} S_i \cdot \left( \sum_{s_i=1}^{S} \sum_{r=1}^{R} \delta_{ir}^{s_i} \right)
    \end{equation}
    
    Целевую функцию можно задать как минимизацию неиспользуемой площади рюкзака:
    
    \begin{equation}
    \min \left( H \cdot W - \sum_{i=1}^{N} S_i \cdot \left( \sum_{s_i=1}^{S} \sum_{r=1}^{R} \delta_{ir}^{s_i} \right) \right)
    \end{equation}
    
    Поскольку величина $H \cdot W$ постоянна для фиксированного рюкзака, обе формулировки эквивалентны и приводят к одному и тому же оптимальному решению. 

	\subsection{Недостатки MIP-модели} %перефразировать, чтобы нормально выглядело как текст
    Несмотря на то, что построенная MIP-модель позволяет точно описать задачу упаковки нерегулярных многоугольников в прямоугольный рюкзак, ее практическое применение связано с рядом ограничений. Преимуществом данной модели является возможность явного учета условий размещения объектов, допустимых поворотов и ограничений на непересечение. Однако точность такой постановки достигается ценой высокой вычислительной сложности.

    При увеличении числа объектов быстро возрастает количество переменных и ограничений модели. Дополнительное усложнение связано с учетом нескольких поворотов каждого объекта, а также с необходимостью описания условий непересечения для пар объектов. В результате даже для наборов средней размерности решение полной MIP-модели требует значительных вычислительных затрат.
    
    Особенно заметно это проявляется в рассматриваемой задаче, где объекты имеют нерегулярную форму. Учет геометрических характеристик таких фигур и использование параметров, основанных на no-fit polygon, повышают точность модели, но одновременно делают ее более громоздкой. Таким образом, вычислительная трудоемкость определяется не только количеством объектов, но и сложностью их геометрического описания.
    
    В практической постановке задачи дополнительным ограничением является необходимость получения решения за разумное время. По этой причине использование полной MIP-модели в качестве единственного инструмента решения оказывается затруднительным. Это делает целесообразным переход к матэвристическому подходу, в котором жадный алгоритм используется для построения начального решения, а математическая модель — для его последующего локального улучшения.
    
    \newpage    
	\section{Жадный алгоритм построения начального решения}
    
    Для построения начального решения в работе используется рандомизированный жадный алгоритм.
    Его идея состоит в том, что фигуры рассматриваются последовательно, однако порядок их
    обработки задаётся случайной перестановкой. Для каждой фигуры рассматриваются все допустимые
    повороты. Для каждого поворота строится множество допустимых положений внутри рюкзака с
    учётом уже размещённых фигур. Из найденных допустимых положений выбирается положение по
    правилу bottom-left, то есть с минимальной координатой $y$, а при равенстве --- с
    минимальной координатой $x$. Если для фигуры существует несколько допустимых поворотов,
    в решение добавляется тот из них, который даёт наилучшее положение по тому же правилу.
    Если ни один поворот не допускает размещения, фигура остаётся неупакованной.

    \begin{algorithm}[tbp]
    \caption{Рандомизированный жадный алгоритм}
    \label{alg:greedy_random_short}
    \small
    \begin{algorithmic}[1]
    \State \textbf{Вход:} множество фигур $\mathcal{I}$, размеры рюкзака $W \times H$, зерно случайности $r$
    \State \textbf{Выход:} начальное жадное решение $Y^g$
    \State Разбить множество фигур на группы поворотов по идентификатору фигуры
    \State Случайным образом задать порядок рассмотрения групп с помощью зерна $r$
    \State $\mathcal{P} \gets \varnothing$

    \ForAll{группа поворотов $G$ в выбранном порядке}
        \State $best \gets \varnothing$
        \ForAll{поворот $i \in G$}
            \State Построить множество допустимых положений $F_i$ фигуры $i$ с учётом уже упакованных фигур $\mathcal{P}$
            \If{$F_i \neq \varnothing$}
                \State Выбрать положение $(x_i, y_i) = \arg \min\limits_{(x,y) \in F_i}(y, x)$
                \If{$best = \varnothing$ или $(x_i, y_i)$ лучше текущего $best$}
                    \State $best \gets (i, x_i, y_i)$
                \EndIf
            \EndIf
        \EndFor
        \If{$best \neq \varnothing$}
            \State Добавить найденное положение в множество размещённых фигур $\mathcal{P}$
        \EndIf
    \EndFor

    \State Сформировать решение $Y^g$ по множеству $\mathcal{P}$
    \State Вычислить значение целевой функции как суммарную площадь упакованных фигур
    \State \Return $Y^g$
    \end{algorithmic}
    \end{algorithm}
    

    \newpage

    \section{Матэвристический алгоритм решения задачи}

    Поскольку рассмотренные выше алгоритмы решения задачи упаковки нерегулярных фигур обладают существенными недостатками, в работе исследуется матэвристический подход, сочетающий построение начального жадного решения и его последующее улучшение с помощью математической модели. При этом локальная MIP-подзадача решается не до оптимальности: достаточно найти допустимое решение, улучшающее исходное жадное размещение.

    Для того, чтобы упростить вычислительную сложность алгоритма, в рамках данной работы было принято решение уменьшить рюкзак, который будет передаваться в решение локальной MIP-подзадачи.

    Как нам известно из пункта 1.6, при большом $N$ решение математической модели является вычислительно трудным. Поэтому, при упаковке с помощью гибридного алгоритма, нам необходимо делать подвыборку предметов среди оставшихся после жадной упаковке в наборе предметов $i = 1, \ldots, N$. 

    В примере, показанном на рисунке, размер $N = 50$ в тестовых данных. Как мы видим, мы упаковали только ААААААА предметов, и у нас осталось $50 - $ААААААА. Так как данное число предметов все еще является слишком большим для алгоритма с поиском допустимого решения математической модели, то нам необходимо сформировать подвыборку предметов. В данном примере это делается случайным образом. 

    \subsection{Модификация математической модели}

    Модель, которую мы приводили в разделе 1, не подходит для добавления в гибридный алгоритм в чистом виде из-за того, что в ней мы упаковываем предметы в неупакованный прямоугольный рюкзак. Данный вариант не подходит, так как нам необходимо учитывать геометрию уже упакованных жадной упаковкой предметов. Поэтому в рамках данной работы была проведена модификация математической модели. 

    Для введения новых ограничений необходимо провести несколько предпросчетов. 

    Первое, необходимо посчитать значение целевой функции локальной подзадачи при расположении объектов жадным алгоритмом. Обозначим данное значение как $S_{greedy}$

    Произведенные изменения: 
      \begin{itemize}
        \item В модель передается новое подмножество предметов $N_{packed}$ - набор уже упакованных объектов. Также передается новые точки отсчетов и их расположение в полу-дискретных координатах $x_{ir}$ и $y_{ir}$, где $i \in N_{packed}$;
        \item Добавляется ограничение на обязательное улучшение решения

        
        $\sum_{i=1}^{N} S_i \cdot \left( \sum_{s_i=1}^{S} \sum_{r=1}^{R} \delta_{ir}^{s_i} \right) > S_{greedy}$;
        \item Поскольку мы формируем новые границы рюкзака, то необходимо поменять ограничения (5) и (6) так, чтобы они учитывали новую границу рюкзака. 
      \end{itemize}

    \newpage

    \subsection{Описание алгоритма}

    Гибридный алгоритм использует жадное решение как начальное, а затем пытается улучшить
    его с помощью локальной математической модели. Сначала строится жадная упаковка. После
    этого в верхней части рюкзака выбираются кандидаты на распаковку. К ним добавляются
    фигуры, которые не были упакованы жадным алгоритмом. Из полученного множества несколько
    раз случайным образом формируется подвыборка фигур, для которой решается локальная
    MIP-подзадача. Остальная часть жадного решения при этом фиксируется. Если найденное
    локальное решение улучшает жадную упаковку по значению целевой функции, то оно
    принимается как новое лучшее решение.

    \begin{algorithm}[tbp]
    \caption{Гибридный алгоритм}
    \label{alg:hybrid_short}
    \small
    \begin{algorithmic}[1]
    \State \textbf{Вход:} множество фигур $\mathcal{I}$, размеры рюкзака $W \times H$, параметры $N_{\text{unpack}}$, $H_{\text{crop}}$, $K_{\text{iter}}$, $K_{\text{sample}}$
    \State \textbf{Выход:} итоговое решение $Y_{\text{final}}$
    \State Построить жадное решение $Y^g$
    \State Выделить верхнюю часть рюкзака высоты $H_{\text{crop}}$
    \State Выбрать множество кандидатов на распаковку $\mathcal{U}$ как последние $N_{\text{unpack}}$ фигур жадного решения, попадающие в выбранную область
    \State Обозначить через $\mathcal{U}^g$ множество фигур, не упакованных жадным алгоритмом
    \State $\Omega \gets \mathcal{U} \cup \mathcal{U}^g$
    \State $Y^{best} \gets Y^g$

    \For{$t \gets 1$ to $K_{\text{iter}}$}
        \State Случайно выбрать подмножество $\Omega_t \subseteq \Omega$ размера не более $K_{\text{sample}}$
        \State Зафиксировать в модели остальную часть жадного решения
        \State Решить локальную MIP-подзадачу для множества $\Omega_t$ в уменьшенном рюкзаке
        \State Получить локальное решение $Y_t$
        \If{$Y_t$ допустимо и $\operatorname{Obj}(Y_t) > \operatorname{Obj}(Y^{best})$}
            \State $Y^{best} \gets Y_t$
        \EndIf
    \EndFor

    \State $Y_{\text{final}} \gets Y^{best}$
    \State \Return $Y_{\text{final}}$
    \end{algorithmic}
    \end{algorithm}

    
    \section{Вычислительные эксперименты}

    из наблюдений (непонятно нужно ли это тут). количество строк влияет на сложность ветвления. количество предметов в случайной выборке влияет на количество нод 
        
    \begin{tcolorbox}[
        colback=red!5,
        colframe=black!50,
        title=Промежуточный план работы,
        sharp corners,
        boxrule=0.5pt,
        left=1mm,right=1mm,top=1mm,bottom=1mm
    ]

    \textbf{до 10.04:}
    \begin{itemize}
        \setlength{\itemsep}{0.15em}
        \setlength{\parskip}{0pt}
        \setlength{\parsep}{0pt}
        \item определить количество строк, которое считается адекватно и не сильно долго
        \item берем "простую обрезку" нового рюкзака
        \begin{itemize}
            \setlength{\itemsep}{0.15em}
            \setlength{\parskip}{0pt}
            \setlength{\parsep}{0pt}
            \item либо просто треть и смотрим, но с ограничением на распаковку
            \item либо нижную границу распакованных предметов х 1.5
            \item макс (2 варианта выше)
        \end{itemize}
        \item случайный выбор для упаковки моделью
        \item получить результаты вычисления для условий выше!!!!
    \end{itemize}
    
    \textbf{до конца апреля:}
    \begin{itemize}
        \setlength{\itemsep}{0.15em}
        \setlength{\parskip}{0pt}
        \setlength{\parsep}{0pt}
        \item конст количество строк
        \item простая обрезка
        \item выбор
        \begin{itemize}
            \setlength{\itemsep}{0.15em}
            \setlength{\parskip}{0pt}
            \setlength{\parsep}{0pt}
            \item случайно (попробовать посчитать как много зависит от случайности)
            \item 10 самых маленьких
            \item 10 самых больших
            \item 10 самых средних
        \end{itemize}
        \item получить результаты вычисления для условий выше!!!!
    \end{itemize}
    
    \end{tcolorbox}

	\newpage
 	
	\newpage
	\section*{ЗАКЛЮЧЕНИЕ}
	\addcontentsline{toc}{section}{ЗАКЛЮЧЕНИЕ}%Не нумеруем заключение, но добавляем в оглавление
	Текст заключения. Еще одна ссылка на литературу \cite{Paper}.
	
	%% БИБЛИОГРАФИЯ
    \newpage
       \begin{thebibliography}{}
    	\addcontentsline{toc}{section}{СПИСОК ЛИТЕРАТУРЫ} %% При таком задании придется повозиться, чтобы задать нумерацию и этому разделу... В тексте она будет, но не в оглавлении
        \bibitem{BLF semi}{Chehrazad S., Roose D., Wauters T.: A fast and scalable bottom-left-fill algorithm to solve nesting problems using a semi-discrete representation. European Journal of Operational Research, \textbf{300}, 809--826 (2022)}
        \bibitem{tutorial}{Bennell J. A., Oliveira J. F.: The geometry of nesting problems: A tutorial. European Journal of Operational Research, \textbf{184}, 397--415 (2008)}
        \bibitem{discr strip}{Umetani S., Murakami S. Coordinate descent heuristics for the irregular strip packing problem of rasterized shapes. European Journal of Operational Research. 2022. Vol. 303. P. 1009--1026.}
        \bibitem{NFP}{Burke E. K., Hellier R. S. R., Kendall G., Whitwell G. Complete and robust no-fit polygon generation for the irregular stock cutting problem. European Journal of Operational Research. 2007. Vol. 179. P. 27--49.}
    	\bibitem{BLF}{Burke E., Hellier R., Kendall G., Whitwell G.: A new bottom-left-fill heuristic algorithm for the two-dimensional irregular packing problem. Operations Research, \textbf{54}, 587--601 (2006)}
    	\bibitem{MIP strip}{Leao A. A. S., Toledo F. M. B., Oliveira J. F., Carravilla M. A.: A semi-continuous MIP model for the irregular strip packing problem. International Journal of Production Research, \textbf{54}(3), 712--721 (2016)}
    \end{thebibliography}
\end{document}
